\section{Introduction}\label{sec:intro}

A central goal of mechanistic interpretability is to understand how neural networks implement learned capabilities at the level of identifiable circuits --- sparse subnetworks of features and connections responsible for specific input--output behaviors \cite{elhage2021mathematical, olsson2022context}. Recent work has made remarkable progress in identifying individual circuits for tasks such as indirect object identification \cite{wang2023interpretability}, greater-than comparison \cite{hanna2023how}, and factual recall \cite{meng2022locating}. Automated methods including activation patching \cite{conmy2023automated}, sparse feature circuits \cite{marks2024sparse}, and attribution-based tracing \cite{templeton2024scaling} have scaled circuit discovery beyond manual analysis. Yet most studies focus on individual circuits for individual tasks, leaving open a fundamental question: are there recurring \emph{structural motifs} --- small subgraph patterns that appear more often than expected by chance --- that characterize how networks organize computation across diverse capabilities?

This question has deep roots in network science. Network motifs, first introduced by \citet{milo2002network} in the analysis of biological and technological networks, provide a principled framework for identifying such recurring patterns. The key idea is to compare the frequency of small subgraph patterns in a real network against their frequency in an ensemble of randomized null models that preserve basic structural properties such as degree sequence. Patterns that appear significantly more often than expected are deemed motifs and are hypothesized to serve as functional building blocks. Among the 13 possible 3-node directed subgraph patterns, the feedforward loop (FFL, also denoted 030T) has emerged as the most functionally significant motif across multiple domains. In biological gene regulatory networks, the FFL consists of a transcription factor that regulates a target gene both directly and indirectly through an intermediate regulator, enabling coherent signal integration, noise filtering, and sign-sensitive delay \cite{alon2007network, mangan2003structure, shenorr2002network}. The ubiquity of FFLs across biological, technological, and social networks \cite{milo2004superfamilies} suggests that this motif captures a fundamental information-processing primitive.

Attribution graphs from the Neuronpedia platform \cite{neuronpedia2024} offer a unique opportunity to test whether similar structural motifs characterize neural network circuits. Each attribution graph is a weighted directed acyclic graph (DAG) where nodes represent sparse autoencoder features \cite{bricken2023monosemanticity} and edges represent attribution scores quantifying causal influence between features for a specific input. Unlike weight-space analyses that examine static network architecture, attribution graphs capture the \emph{dynamic} computation path for individual inputs, making them ideal for studying input-specific circuit structure. Crucially, Neuronpedia provides attribution graphs across multiple capability domains --- from antonym detection and arithmetic to code completion and multi-hop reasoning --- enabling cross-domain structural comparisons at unprecedented scale.

In this work, we systematically analyze FFL motifs across a large corpus of 200 Neuronpedia attribution graphs spanning 8 capability domains. Our analysis addresses five hypotheses that together characterize the role of FFL motifs as organizational primitives:

\begin{enumerate}
\item \textbf{H1 (FFL Universality):} FFL motifs are statistically over-represented across all capability domains, not just specific task types. This would suggest that FFLs reflect a general computational strategy rather than a task-specific artifact.
\item \textbf{H2 (Hub Structural Importance):} Nodes participating in many FFL instances (``hubs'') are disproportionately important for circuit function, as measured by the downstream impact of node ablation. This would establish a causal link between motif participation and functional significance.
\item \textbf{H3 (Weighted Feature Discrimination):} Weighted motif features that capture edge-weight information (intensity, sign coherence, path dominance) discriminate capability domains better than binary motif counts alone. This would demonstrate that the quantitative properties of motifs, not just their topology, carry domain-relevant information.
\item \textbf{H4 (Unique Information):} Motif-level features carry structural information about capability domains that is genuinely orthogonal to --- not redundant with --- graph-level statistics such as node count, edge count, density, and degree distribution.
\item \textbf{H5 (Higher-Order Derivativeness):} Universal 4-node motifs are fully embedded within FFL structures, indicating that the 3-node FFL is the fundamental building block from which higher-order motif structure emerges.
\end{enumerate}

The network motif framework offers a natural vocabulary for this investigation. Network motifs --- small subgraph patterns that recur significantly more often than expected in randomized null models --- have proven invaluable for understanding functional architecture in biological, technological, and social networks \cite{milo2002network, milo2004superfamilies}. Among the 13 possible 3-node directed subgraph patterns, the feedforward loop (FFL, also denoted 030T) has received the most attention due to its functional versatility: in gene regulatory networks, the coherent FFL acts as a persistence detector while the incoherent FFL generates adaptive pulses \cite{alon2007network}. If analogous motifs characterize neural network attribution circuits, it would suggest deep parallels between biological and artificial information-processing architectures. Moreover, identifying such motifs would provide a principled decomposition of circuit structure into interpretable building blocks, complementing existing approaches that focus on individual features or complete circuits.

The practical stakes of this investigation extend beyond scientific curiosity. If universal structural motifs exist in neural network circuits, they could serve as (i) structured starting points for circuit discovery, reducing the search space from exponentially many subgraphs to a small set of canonical patterns; (ii) compact fingerprints for comparing how different models implement similar capabilities; (iii) anomaly detectors for flagging unusual or potentially unreliable circuit implementations; and (iv) targets for architecture-aware pruning that selectively preserves functionally critical circuit elements. These applications motivate a thorough investigation of motif structure in attribution circuits.

Our contributions are as follows. First, we establish FFL universality across all 8 capability domains using a rigorous null-model framework with 50 degree-preserving random graphs per attribution graph, supplemented by Benjamini--Hochberg FDR correction, pruning robustness analysis, convergence diagnostics, and layer-preserving null validation. We report median FFL Z-scores of 47.2 across all domains (\S\ref{sec:results_h1}). Second, we quantify hub structural importance via graph-theoretic node ablation on 11,185,566 enumerated FFL instances, finding that hub removal causes 1.92$\times$ greater downstream attribution loss than layer-matched controls, with a strong dose--response relationship (\S\ref{sec:results_h2}). Third, we show that weighted motif features achieve spectral clustering $NMI=0.705$ against domain labels, far exceeding binary motif counts at $NMI=0.101$, with combined features reaching $NMI=0.844$ (\S\ref{sec:results_h3}). Fourth, we demonstrate unique motif information via McFadden variance decomposition (unique $R^2=0.0184$, bootstrap CI excludes 0) and residualized clustering (\S\ref{sec:results_h4}). Fifth, we show that all 4 universal 4-node motif types exhibit 100\% FFL containment, confirming that higher-order structure is entirely FFL-derivative (\S\ref{sec:results_h5}).

Together, these results provide the first comprehensive evidence that FFL motifs serve as fundamental organizational primitives of neural network attribution circuits. Our findings have implications for circuit discovery, model comparison, anomaly detection, and the design of interpretability tools. We discuss these implications alongside limitations and future directions in \S\ref{sec:discussion}. All data, code, and statistical analyses are provided for full reproducibility.

The remainder of this paper is organized as follows. Section~\ref{sec:related} reviews related work on network motifs, mechanistic interpretability, and weighted graph analysis. Section~\ref{sec:methods} describes our data collection pipeline, motif enumeration methodology, ablation framework, weighted feature extraction, and unique information decomposition approach. Sections~\ref{sec:results_h1}--\ref{sec:results_h5} present results for each of the five hypotheses, including detailed statistical analyses and honest assessments of negative or ambiguous findings. Section~\ref{sec:discussion} discusses implications, limitations, and future directions, and Section~\ref{sec:conclusion} concludes. Throughout, we emphasize transparent reporting of both positive and negative results, including methodological caveats regarding FDR correction limitations and the modest absolute magnitude of unique motif information. We believe this transparency strengthens rather than weakens the overall contribution.